\chapter{Исчисление высказываний}%
%\index{Исчисление высказываний}
        \label{propositional-calculus}%
Напомним, что тавтологии%
\index{Тавтология}
\т это пропозициональные формулы,%
\index{Пропозициональная формула}
истинные при всех значениях переменных. Оказывается, что все
тавтологии можно получить из некоторого набора \лк аксиом\пк\ с
помощью \лк правил вывода\пк, которые имеют чисто синтаксический
характер и никак не апеллируют к смыслу формулы, её истинности
\итд Эту задачу решает так называемое \emph{исчисление
высказываний}.%
\index{Исчисление!высказываний|defin}
В этой главе мы перечислим аксиомы и правила
вывода этого исчисления, и приведём несколько доказательств
\emph{теоремы о полноте}%
\index{Теорема!о полноте ИВ}
(которая утверждает, что всякая
тавтология%
\index{Тавтология}
выводима в исчислении выс\-казываний).

\section{Исчисление высказываний (ИВ)}
        \label{propositional-calculus-axioms}%

Каковы бы ни были формулы $A,B,C$, следующие формулы называют
\emph{аксиомами исчисления высказываний}:%
\index{Аксиомы!исчисления высказываний|defin}%
\index{Схемы аксиом исчисления высказываний|defin}%
\index{Исчисление!высказываний!схемы аксиом|defin}

\myboxrighttxt{(100)}{(1)} $A\to(B\to A)$;

\myboxrighttxt{(100)}{(2)} $(A\to(B\to C))\to((A\to B)\to
(A \to C))$;

\myboxrighttxt{(100)}{(3)} ($A\land B)\to A$;

\myboxrighttxt{(100)}{(4)} ($A\land B)\to B$;

\myboxrighttxt{(100)}{(5)} $A\to(B\to(A\land B))$;

\myboxrighttxt{(100)}{(6)} $A\to (A\lor B)$;

\myboxrighttxt{(100)}{(7)} $B\to(A\lor B)$;

\myboxrighttxt{(100)}{(8)} $(A\to C) \to ( (B\to C) \to
((A\lor B)\to C))$;

\myboxrighttxt{(100)}{(9)} $\neg A\to(A\to B)$;

\myboxrighttxt{(100)}{(10)} $(A\to B)\to
((A\to \lnot B)\to\lnot A)$;

\myboxrighttxt{(100)}{(11)} $A\lor \neg A$.


\noindent
Как говорят, мы имеем здесь одиннадцать \лк схем аксиом\пк; из каждой
схемы можно получить различные конкретные аксиомы, заменяя
входящие в неё буквы на пропозициональные формулы.

Единственным правилом вывода исчисления высказываний является
правило со средневековым названием \лк modus ponens\пк\ (MP).%
\index{Modus ponens (MP), правило|defin}%
\index{Правило!modus ponens (MP)|defin}
Это правило разрешает получить (вывести) из формул $A$ и $(A\to
B)$ формулу~$B$.

\emph{Выводом}%
\index{Вывод в исчислении высказываний|defin}%
\index{Исчисление!высказываний!вывод|defin}
в исчислении высказываний называется
конечная последовательность формул, каждая из которых есть
аксиома или получается из предыдущих по правилу modus ponens.

Вот пример вывода (в нём первая формула является частным случаем
схемы~(1), вторая\т схемы~(2), а последняя получается из
двух предыдущих по правилу modus ponens):
\begin{align*}
   &(p\to(q\to p)),\\
   &(p\to(q\to p))\to((p\to q)\to(p\to p)),\\
   &((p\to q)\to(p\to p)).
\end{align*}

Пропозициональная формула $A$ называется \emph{выводимой}%
\index{Формула!выводимая в исчислении высказываний|defin}%
\index{Выводимая формула|defin}%
\index{Исчисление!высказываний!выводимая формула|defin}
в исчислении высказываний, или \emph{теоремой}%
\index{Теорема!исчисления высказываний|defin}%
\index{Исчисление!высказываний!теорема|defin}
исчисления
высказываний, если существует вывод, в котором последняя формула
равна $A$. Такой вывод называют выводом формулы $A$. (В принципе
можно было бы и не требовать, чтобы формула $A$ была последней\т
все дальнейшие формулы можно просто вычеркнуть.)

Как мы уже говорили, в исчислении высказываний выводятся все
тавтологии и только они. Обычно это утверждение разбивают на две
части: простую и сложную. Начнём с простой:

\begin{theorem}[о корректности ИВ]
\index{Теорема!о корректности ИВ|defin}%
    Всякая теорема исчисления высказываний есть тавтология.
\end{theorem}

\begin{proof}
        %
Несложно проверить, что все аксиомы\т тавтологии. Для примера
проделаем это для самой длинной аксиомы (точнее, схемы аксиом)\т
для второй. В каком случае формула
        $$
(A\to(B\to C))\to((A\to B)\to (A\to C))
        $$
(где $A,B,C$\т некоторые формулы) могла бы быть ложной? Для
этого посылка $A\to(B\to C)$ должна быть истинной, а заключение
$(A\to B)\to (A\to C)$\т ложным. Чтобы заключение было ложным,
формула $A\to B$ должна быть истинной, а формула $A\to C$\т
ложной. Последнее означает, что $A$ истинна, а $C$ ложна. Таким
образом, мы знаем, что $A$, $(A\to B)$ и $(A\to(B\to C))$
истинны. Отсюда следует, что $B$ и $(B\to C)$ истинны,
и потому $C$
истинна\т противоречие. Значит, наша формула не бывает ложной.

Корректность правила MP также ясна: если формулы $(A\to B)$
и $A$ всегда истинны, то по определению импликации формула~$B$
также всегда истинна. Таким образом, все формулы, входящие в
выводы (все теоремы) являются тавтологиями.
        %
\end{proof}

Гораздо сложнее доказать обратное утверждение.

\begin{theorem}[о полноте ИВ]
\index{Теорема!о полноте ИВ|defin}%
        \label{propositional-completeness}%
Всякая тавтология есть теорема исчисления высказываний.
        %
\end{theorem}

\begin{proof}
        %
Мы предложим несколько альтернативных доказательств этой
теоремы. Но прежде всего мы должны приобрести некоторый опыт
построения выводов и использования аксиом.

\textsf{Лемма 1}. Какова бы ни была формула $D$, формула ${(D\to
D)}$ является теоремой.

Докажем лемму, предъявив вывод формулы $(D\to D)$ в исчислении
высказываний.

{
\hangindent=4mm
\hangafter=1
\rightskip=0pt plus1fil
\noindent
1.~$({D\to((D\to D)\to D))}\hm\to{((D\to(D \to D))\to(D\to D))}$\\\relax
[аксиома~2 при $A=D$, $B=(D\to D)$, $C=D$];

\raggedright
2. $D\to((D\to D)\to D)$ [аксиома~1];

3. $(D\to(D\to D))\to(D\to D)$ [из~1 и~2 по правилу~MP];

4. $D\to(D\to D)$ [аксиома~1];

5. $(D\to D)$ [из~3 и~4 по правилу MP].
}

\medskip

Как видно, вывод даже такой простой тавтологии, как $(D\to D)$,
требует некоторой изобретательности. Мы облегчим себе жизнь,
доказав некоторое общее утверждение о выводимости.

Часто мы рассуждаем так: предполагаем, что выполнено какое\д то
утверждение $A$, и выводим различные следствия. После того как
другое утверждение~$B$ доказано, мы вспоминаем, что использовали
предположение~$A$, и заключаем, что мы доказали утверждение
$A\to B$. Следующая лемма, называемая иногда \лк леммой о
дедукции\пк, показывает, что этот подход правомерен и для
исчисления высказываний.

Пусть $\Gamma$\т некоторое множество формул. \emph{Выводом из
$\Gamma$}%
\index{Вывод из $\Gamma$|defin}
называется конечная последовательность формул, каждая
из которых является аксиомой, принадлежит $\Gamma$ или
получается из предыдущих по правилу MP. (Другими словами, мы как
бы добавляем формулы из $\Gamma$ к аксиомам исчисления
высказываний\т именно как формулы, а не как схемы аксиом.)
Формула~$A$ \emph{выводима из}~$\Gamma$,%
\index{Формула!выводимая из $\Gamma$|defin}
если существует вывод
из $\Gamma$, в котором она является последней формулой. В этом
случае мы пишем $\Gamma\vdash A$. Если $\Gamma$ пусто, то речь
идёт о выводимости в исчислении высказываний, и вместо
$\varnothing\vdash A$ пишут просто $\vdash A$.

        \label{deduction-lemma}%
\textsf{Лемма 2 (о дедукции)}.%
\index{Лемма!о дедукции|defin}
$\Gamma\vdash A\to B$ тогда и только тогда, когда
$\Gamma\cup\{A\}\vdash B$.

В одну сторону утверждение почти очевидно: пусть $\Gamma\hm\vdash
(A\to B)$. Тогда и $\Gamma, A\vdash(A\to B)$. (Для краткости мы
опускаем фигурные скобки и заменяем знак объединения запятой.)
Согласно определению, $\Gamma, A\vdash A$, откуда по MP получаем $\Gamma,
A\vdash B$.

Пусть теперь $\Gamma,A\vdash B$. Нам надо построить вывод
формулы $A\to B$ из $\Gamma$. Возьмём вывод $C_1,C_2,\ldots,C_n$
формулы $B=C_n$ из $\Gamma, A$. Припишем ко всем формулам этого
вывода слева посылку $A$:
        $$
(A\to C_1), (A\to C_2),\dots,(A\to C_n).
        $$
Эта последовательность оканчивается на $(A\to B)$. Сама по себе
она не будет выводом из $\Gamma$, но из неё можно получить такой
вывод, добавив недостающие формулы, и тем самым доказать лемму о
дедукции.

Будем добавлять эти формулы, двигаясь слева направо. Пусть мы
подошли к формуле $(A\to C_i)$. По предположению формула $C_i$
либо совпадает с $A$, либо принадлежит $\Gamma$, либо является
аксиомой, либо получается из двух предыдущих по правилу MP.
Рассмотрим все эти случаи по очереди.

(1) Если $C_i$ есть $A$, то очередная формула имеет вид $(A\hm\to
A)$. По лемме~1 она выводима, так что перед ней мы добавляем её
вывод.

(2) Пусть $C_i$ принадлежит $\Gamma$. Тогда мы вставляем формулы
$C_i$ и $C_i\to(A\to C_i)$ (аксиома~1). Применение правила MP к
этим формулам даёт $(A\to C_i)$, что и требовалось.

(3) Те же формулы можно добавить, если $C_i$ является аксиомой
исчисления высказываний.

(4) Пусть, наконец, формула
$C_i$ получается из двух предыдущих формул
по правилу MP. Это значит, что в исходном выводе ей предшествовали
формулы
$C_j$ и $(C_j\hm\to C_i)$. Тогда в новой последовательности (с
добавленной посылкой~$A$) уже были формулы $(A\to C_j)$ и
$(A\to(C_j\hm\to C_i))$. Поэтому мы можем продолжить наш $\Gamma$\д
вывод, написав формулы

        \smallskip
\noindent
$((A\to(C_j\to C_i))\to((A\to C_j)\to (A\to C_i))$ (аксиома 2);

\noindent
$((A\to C_j)\to (A\to C_i))$ (modus ponens);

\noindent
$(A\to C_i)$ (modus ponens).
        \smallskip

Итак, во всех четырёх случаях мы научились дополнять
последовательность до вывода из $\Gamma$, так что лемма о
дедукции доказана.

\begin{problem}
        %
Докажите, что для любых формул $A,B,C$ формула
        $$
(A\to B)\to((B\to C)\to (A\to C))
        $$
выводима в исчислении высказываний. (Указание:
используйте лемму о дедукции и тот факт, что ${A\to B}, {B\to C}, A
\hm\vdash C$.)
        %
\end{problem}

\begin{problem}
        %
Докажите, что если $\Gamma_1\vdash A$ и $\Gamma_2,A\vdash B$, то
$\Gamma_1\hm\cup\Gamma_2\hm\vdash B$. (Это свойство иногда называют
\лк правилом сечения\пк\ (cut);%
   \index{Правило!сечения|defin}%
   \index{Сечения!правило|defin}
говорят, что формула $A$
\лк отсекается\пк\ или \лк высекается\пк.
Сходные правила играют центральную роль в теории
доказательств, где формулируется и доказывается \лк теорема об
устранении сечения\пк\ для различных логических систем.)
        %
\end{problem}

\begin{problem}
        %
Добавим к исчислению высказываний, помимо правила MP,
ещё одно правило, называемое \emph{правилом подстановки}.%
\index{Правило!подстановки|defin}%
\index{Подстановки правило|defin}
Оно разрешает заменить в выведенной формуле все переменные на
произвольные формулы (естественно, вхождения одной переменной
должны заменяться на одну и ту же формулу). Покажите, что после
добавления такого правила класс выводимых формул не изменится,
но лемма о дедукции перестанет быть верной.
        %
\end{problem}

Заметим, что мы пока что использовали только две первые аксиомы
исчисления высказываний. Видно, кстати, что они специально
подобраны так, чтобы прошло доказательство леммы о дедукции.

Другие аксиомы описывают свойства логических связок. Аксиомы $3$
и $4$ говорят, какие следствия можно вывести из конъюнкции
($A\land B \vdash A$ и $A\land B \vdash B$). Напротив, аксиома~5
говорит, как можно вывести конъюнкцию. Из неё легко следует
такое правило: если $\Gamma\vdash A$ и $\Gamma\vdash B$, то
$\Gamma\vdash(A\land B)$ (применяем эту аксиому и дважды правило~MP).
Часто подобные правила записывают так:
        $$
\frac{\Gamma\vdash A \qquad \Gamma\vdash B}{\Gamma\vdash A\land B}
        $$
(над чертой пишут \лк посылки\пк\ правила, а снизу\т его \лк
заключение\пк, вытекающее из посылок).

\begin{problem}
        %
Докажите, что формула $(A\hm\to{(B\to C)})\hm\to ({(A\land
B)}\hm\to C)$, так же как и обратная к ней формула (в которой
посылка и заключение переставлены), являются теоремами
исчисления высказываний. Докажите аналогичное утверждение про
формулы ${(A\land B)}\hm\to{(B\land A)}$ и $((A\land B)\land C)\hm\to
(A\hm\land{(B\land C)})$.
        %
\end{problem}

Аксиомы 6\ч 7 позволяют утверждать, что $A\vdash A\lor B$ и
$B\vdash A\lor B$. Аксиома~8 обеспечивает такое правило:
        $$
\frac{\Gamma, A \vdash C \qquad \Gamma,B\vdash C}{\Gamma, A\lor B\vdash C}
        $$
Оно соответствует такой схеме рассуждения: \лк Предположим, что
$A\lor B$. Разберём два случая. Если выполнено $A$, то
\hole и потому~$C$. Если выполнено $B$,
то~\hole и потому~$C$. В обоих случаях верно~$C$.
Значит,~$A\lor B$ влечёт~$C$.\пк

Обоснование: дважды воспользуемся
леммой о дедукции, получив $\Gamma\hm\vdash (A\hm\to C)$ и
$\Gamma\hm\vdash(B\hm\to C)$, а затем дважды применим правило~MP
к этим формулам и аксиоме ${(A\to C)}\hm\to ({(B\to C)}\hm\to
({(A\lor B)}\hm\to C))$.
Получив формулу ${(A\lor B)}\hm\to C$, опять применим правило~MP
к ней и формуле $(A\lor B)$.

\begin{problem}
Докажите, что следующие формулы, а также обратные к ним (меняем
местами посылку и заключение) являются теоремами исчисления
высказываний:
        \begin{align*}
((A\lor B)\to C)&\to ((A\to C)\land(B\to C)),\\
((A\land C)\lor (B\land C))&\to ((A\lor B)\land C),\\
((A\lor C)\land (B\lor C))&\to ((A\land B)\lor C).
        \end{align*}
\end{problem}

У нас остались ещё три аксиомы, касающиеся
отрицания. Аксиома~9 гарантирует, что из противоречивого набора
посылок можно вывести что угодно: если $\Gamma\vdash A$ и
$\Gamma\vdash\lnot A$, то $\Gamma\vdash B$ для любого $B$.
Аксиома 10, напротив, объясняет, как можно вывести отрицание
некоторой формулы
$A$: надо допустить $A$ и вывести два противоположных заключения
$B$ и $\lnot B$. Точнее говоря, имеет место такое правило:
        $$
\frac{\Gamma, A\vdash B\qquad \Gamma, A\vdash\lnot B}
     {\Gamma \vdash \lnot A}
        $$
(в самом деле, дважды применяем лемму о дедукции, а затем
правило MP с аксиомой~10).

Аксиомы 9 и 10 позволяют вывести некоторые логические законы,
связанные с отрицанием. Докажем, например, что (для любых формул
$A$ и $B$) формула
        $$
(A\to B)\to(\lnot B\to\lnot A)
        $$
(\лк закон контрапозиции\пк)
является теоремой исчисления высказываний. В самом деле, по
лемме о дедукции достаточно установить, что
        $$
(A\to B), \lnot B \vdash \lnot A.
        $$
Для этого, в свою очередь, достаточно вывести из трёх посылок $(A\hm\to B),
\lnot B, A$ какую\д либо формулу и её отрицание (в данном случае
формулы~$B$ и~$\lnot B$).

\begin{problem}
        %
Выведите формулы $A\to\lnot\lnot A$ и $\lnot\lnot\lnot A\to\lnot
A$ с помощью аналогичных рассуждений.
        %
\end{problem}

Последняя аксиома, называемая \лк законом исключённого
третьего\пк,%
\index{Закон!исключённого третьего|defin}%
\index{Исключённого третьего закон|defin}
и иногда читаемая как \лк третьего не дано\пк\
(tertium
  \index{tertium non datur, закон|defin}%
non datur в латинском оригинале), вызвала в первой
половине века большое количество споров.
(См.~раздел~\ref{intuitionism} об интуиционистской логике,%
\index{Интуиционистская логика}
в которой этой аксиомы нет.)

Из неё можно вывести закон \лк снятия двойного отрицания\пк,%
  \index{Закон!снятия двойного отрицания|defin}%
  \index{Двойное отрицание, снятие|defin}
формулу $\lnot\lnot A\to A$. Достаточно
показать, что $A\lor \lnot A, \lnot\lnot A \vdash A$. По правилу
разбора случаев, достаточно установить, что $A, \lnot\lnot A
\vdash A$ (это очевидно) и что $\lnot A, \lnot\lnot A \vdash
A$ (а это верно, так как из двух противоречащих друг другу
формул выводится что угодно с помощью аксиомы~9).

\begin{problem}
        %
Докажите, что формула $(\lnot B\to\lnot A)\to(A\to B)$ является
теоремой исчисления высказываний. (Указание: используйте закон
исключённого третьего.)
        %
\end{problem}

\begin{problem}
        %
Исключим из числа аксиом исчисления высказываний закон
исключённого третьего, заменив его на закон снятия двойного
отрицания. Покажите, что от этого класс выводимых формул не
изменится.
        %
\end{problem}

\begin{problem}
        %
Докажите, что при наличии аксиомы исключённого третьего (11)
аксиома (10) является лишней\т её (точнее следовало бы сказать:
любой частный случай этой схемы аксиом) можно вывести из
остальных аксиом.
        %
\end{problem}

Теперь уже можно доказать теорему о полноте: всякая тавтология
выводима в исчислении высказываний. Идея доказательства состоит
в разборе случаев. Поясним её на примере. Пусть $A$\т
произвольная формула, содержащая переменные $p,q,r$.
Предположим, что $A$ истинна, когда все три переменные истинны.
Тогда, как мы докажем,
        $
   p,q,r \vdash A.
        $
Вообще каждой строке таблицы истинности для формулы $A$
соответствует утверждение о выводимости. Например, если $A$
ложна, когда $p$ и $q$ ложны, а $r$ истинно, то
        $
\lnot p, \lnot q, r \vdash \lnot A.
        $
Если формула $A$ является тавтологией, то окажется, что она
выводима из всех восьми возможных вариантов посылок. Пользуясь
законом исключённого третьего, можно постепенно избавляться от
посылок. Например, из $p,q,r\hm\vdash A$ и $p,q,\lnot r\vdash A$
можно получить $p,q,(r\lor\lnot r)\vdash A$, то есть $p,q\vdash
A$ (поскольку $(r\lor\lnot r)$ является аксиомой).

Проведём это рассуждение подробно. Начнём с такой леммы:

\textsf{Лемма 3}. Для произвольных формул $P$ и $Q$
\begin{align*}
P,Q&\vdash (P\land Q);& P,Q&\vdash (P\lor Q);\\
P,\lnot Q&\vdash \lnot (P\land Q);& P,\lnot Q&\vdash (P\lor Q);\\
\lnot P,Q&\vdash \lnot (P\land Q);& \lnot P,Q&\vdash (P\lor Q);\\
\lnot P, \lnot Q;&\vdash \lnot (P\land Q)& \lnot P,
\lnot Q&\vdash \lnot (P\lor Q);\\[1.5ex]
P,Q&\vdash (P\to Q);&&\\
P,\lnot Q&\vdash \lnot (P\to Q);& P&\vdash \lnot (\lnot P);\\
\lnot P,Q&\vdash (P\to Q);& \lnot P&\vdash \lnot P.\\
\lnot P, \lnot Q&\vdash (P\to Q);&&
        %
\end{align*}

Эта лемма говорит, что если принять в качестве гипотез
истинность или ложность формул $P$ и $Q$, являющихся частями
конъюнкции, дизъюнкции или импликации, то можно будет доказать
или опровергнуть всю формулу (в зависимости от того, истинна она
или ложна). Последняя часть содержит аналогичное утверждение про
от\-рицание.

После предпринятой нами тренировки доказать эти утверждения
несложно. Например, убедимся, что $\lnot P\hm\vdash\lnot(P\land
Q)$. Для этого достаточно вывести два противоположных
утверждения из $\lnot P, (P\land Q)$; ими будут утверждения $P$
и $\lnot P$.

Проверим ещё одно утверждение: $\lnot P, \lnot Q\vdash\lnot
(P\lor Q)$. Нам надо вывести два противоположных утверждения из
$\lnot P, \lnot Q, (P\lor Q)$. Покажем, что из $\lnot P, \lnot
Q, (P\lor Q)$ следует всё, что угодно. По правилу разбора
случаев достаточно убедиться, что из $\lnot P, \lnot Q, P$ и из
$\lnot P, \lnot Q, Q$ следует всё, что угодно\т но это мы знаем.

Утверждения, касающиеся импликации, просты: в самом деле, мы
знаем, что $Q\vdash (P\to Q)$ благодаря аксиоме~1, а $\lnot
P\vdash (P\to Q)$ благодаря аксиоме~9.

Остальные утверждения леммы столь же просты.

Теперь мы можем сформулировать утверждение о разборе случаев для
произвольной формулы.

\textsf{Лемма 4}. Пусть $A$\т произвольная формула, составленная
из переменных $p_1,\ldots,p_n$. Тогда для каж\-дой строки таблицы
истинности формулы~$A$ имеет место соответствующее утверждение о
выводимости: если $\varepsilon_1,\dots,\varepsilon_n,
\varepsilon\hm\in\{0,1\}$, и значение формулы $A$ есть
$\varepsilon$ при $p_1=\varepsilon_1,\dots,p_n\hm=\varepsilon_n$,
то
        $$
\lnot_{\varepsilon_1} p_1,\dots,\lnot_{\varepsilon_n} p_n
        \vdash
\lnot_{\varepsilon A},
        $$
где $\lnot_u \varphi$ обозначает $\varphi$ при $u=1$ и
$\lnot \varphi$ при $u=0$ (напомним, что $1$ обозначает истину,
а $0$\т ложь).

Лемма очевидно доказывается индукцией по построению формулы $A$.
Мы имеем посылки, утверждающие истинность или ложность
переменных, и для всех подформул (начиная с переменных и идя ко
всей формуле) выводим их или их отрицания с помощью леммы~3.

Если формула $A$ является тавтологией, то из всех
$2^n$ вариантов посылок выводится именно она, а не её отрицание.
Тогда правило разбора случаев и закон исключённого третьего
позволяют избавиться от посылок: сгруппируем их в пары,
отличающиеся в позиции $p_1$ (в одном наборе посылок стоит
$p_1$, в другом $\lnot p_1$), по правилу разбора случаев заменим
их на посылку $(p_1\lor\lnot p_1)$, которую можно выбросить (она является
аксиомой). Сделав так для всех пар, получим
$2^{n-1}$ выводов, в посылках которых нет $p_1$; повторим этот
процесс с посылками $p_2$, $\lnot p_2$
\итд В конце концов мы убедимся, что формула $A$
выводима без посылок, как и утверждает
теорема о полноте.
        %
\end{proof}

\section{Второе доказательство теоремы о полноте}
        \label{completeness-second-proof}%

Это доказательство, в отличие от предыдущего, обобщается
на более сложные случаи (исчисление предикатов, интуиционистское исчисление
высказываний).%
\index{Исчисление!предикатов}

Начнём с такого определения: множество формул $\Gamma$
называется \emph{совместным},%
\index{Совместное множество формул|defin}%
\index{Множество формул!совместное|defin}
если существует набор значений
переменных, при которых все формулы из $\Gamma$ истинны.
Заметим, что формула $\varphi$ является тавтологией тогда и
только тогда, когда множество, состоящее из единственной формулы
$\lnot \varphi$, не является совместным. Для случая одной формулы есть
специальный термин: формула $\tau$ \emph{выполнима},%
\index{Выполнимая формула|defin}%
\index{Формула!выполнимая|defin}
если
существуют значения переменных, при которых она истинна, то есть
если множество $\{\tau\}$ совместно. Тавтологии\т это формулы,
отрицания которых не выполнимы.

Множество формул $\Gamma$ называется \emph{противоречивым},%
\index{Противоречивое множество формул|defin}%
\index{Множество формул!противоречивое|defin}
если
из него одновременно выводятся формулы $A$ и $\lnot A$. Мы
знаем, что в этом случае из него выводятся вообще все формулы.
(В противном случае $\Gamma$ называется
\emph{непротиворечивым}.)%
\index{Непротиворечивое множество формул|defin}%
\index{Множество формул!непротиворечивое|defin}

\begin{theorem}[корректность исчисления высказываний, вторая форма]
\index{Теорема!о корректности ИВ, вторая форма|defin}%
        \label{correctness-second-form}%
Всякое совместное множество формул непротиворечиво.
        %
\end{theorem}

\begin{proof}
        %
В самом деле, пусть совместное множество $\Gamma$ противоречиво.
Так как оно совместно, существуют значения переменных, при которых
все формулы из $\Gamma$ истинны. С~другой стороны, из $\Gamma$
выводится некоторая формула~$B$ и её отрицание. Может ли так быть?

Оказывается, что нет. Мы уже видели, что всякая выводимая
формула истинна при всех значениях переменных (является
тавтологией). Справедливо и несколько более общее утверждение:
если $\Gamma\vdash A$ и при некоторых значениях переменных
все формулы из
$\Gamma$ истинны, то и формула $A$ истинна при этих значениях
переменных. (Как и раньше, это легко доказывается индукцией по
построению вывода $A$ из $\Gamma$.)

В нашей ситуации это приводит к тому, что на выполняющем наборе
значений переменных для $\Gamma$ должны быть истинны обе формулы
$B$ и $\lnot B$, что, разумеется, невозможно.
        %
\end{proof}

Мы называем это утверждение другой формой теоремы о корректности
исчисления высказываний, поскольку из него формально можно
вывести, что всякая теорема является тавтологией: если $A$\т
теорема, то множество~$\{\lnot A\}$ противоречиво (из него
выводятся $A$ и $\lnot A$), потому несовместно, значит, $\lnot
A$ всегда ложна, поэтому $A$ всегда истинна.

\begin{theorem}[полнота исчисления высказываний, вторая форма]
\index{Теорема!о полноте ИВ, вторая форма|defin}%
        \label{completeness-second-form}%
Всякое непротиворечивое множество сов\-местно.%
        %
\end{theorem}

\begin{proof}
        %
Нам дано непротиворечивое множество $\Gamma$, а надо найти такие
значения переменных, при которых все формулы из $\Gamma$
истинны. (Вообще говоря, множество $\Gamma$ может быть
бесконечно и содержать бесконечное число разных переменных.)

Пусть есть какая\д то переменная $p$, встречающаяся
в формулах из семейства $\Gamma$. Нам надо решить, сделать ли
её истинной или ложной. Если оказалось так, что из
$\Gamma$ выводится формула $p$, то выбора нет: она обязана
быть истинной в тех наборах, где формулы из $\Gamma$ истинны
(как мы видели при доказательстве корректности). По тем же
причинам, если из $\Gamma$ выводится $\lnot p$, то в выполняющем
наборе переменная $p$ обязательно будет ложной.

Если оказалось так, что для любой переменной $p$ либо она сама,
либо её отрицание выводятся из $\Gamma$, то выполняющий набор
значений определён однозначно, и надо только проверить, что он
действительно будет выполняющим. А если для каких\д то
переменных нельзя вывести ни их, ни их отрицание, то мы пополним
наш набор $\Gamma$ так, чтобы они, как теперь модно говорить,
\лк определились\пк.

Проведём это рассуждение подробно.
Рассмотрим все переменные, входящие в какие\д либо формулы из
множества $\Gamma$; обозначим множество этих переменных
через~$V$. Зафиксируем это множество и до конца доказательства
теоремы о полноте будем рассматривать только формулы с переменными
из множества~$V$, не оговаривая этого~особо.

Назовём непротиворечивое множество $\Gamma$ \emph{полным},%
\index{Полное множество формул|defin}%
\index{Множество формул!полное|defin}
если
для любой формулы $F$ имеет место либо $\Gamma\vdash F$, либо
$\Gamma\vdash \lnot F$ (одновременно этого быть не может, так
как $\Gamma$~непротиворечиво).

Утверждение теоремы о полноте очевидно следует из двух лемм:

        \label{completeness-second-proof-lemma-1}%
\textsf{Лемма 1.} Всякое непротиворечивое множество $\Gamma$
содержится в непротиворечивом полном множестве $\Delta$.

        \label{completeness-second-proof-lemma-2}%
\textsf{Лемма 2.} Для всякого непротиворечивого полного
множества $\Delta$ существует набор значений переменных
(из~$V$, напомним), при котором все формулы из~$\Delta$
истинны.

Доказательство леммы 1. Основную роль здесь играет такое
ут\-верж\-де\-ние: если $\Gamma$\т непротиворечивое множество, а $A$\т
произвольная формула, то хотя бы одно из множеств
$\Gamma\cup\{A\}$ и $\Gamma\cup\{\lnot A\}$ непротиворечиво. В
самом деле, если оба множества
$\Gamma\cup\{A\}$ и
$\Gamma\cup\{\lnot A\}$ противоречивы, то
$\Gamma\vdash\lnot A$ и $\Gamma\vdash\lnot\lnot A$,
но множество $\Gamma$
предполагалось непротиворечивым.

Если множество переменных $V$ конечно или счётно, то
доказательство леммы~$1$ легко завершить: множество всех формул
тогда счётно, и просматривая их по очереди, мы можем добавлять к
$\Gamma$ либо саму формулу, либо её отрицание, сохраняя
непротиворечивость. Получится, очевидно, полное множество.
Чуть менее очевидна его непротиворечивость: оно было
непротиворечиво на каждом шаге, но почему предельное
множество (объединение возрастающей последовательности)
будет непротиворечиво? Дело в том, что в выводе двух
противоречащих друг другу формул может быть задействовано
только конечное число формул из $\Gamma$ (по определению
выводимости: вывод есть конечная последовательность формул).
Поэтому все эти формулы должны появиться на некотором
конечном шаге конструкции, а это невозможно (на всех
шагах множество непротиворечиво).

Для случая произвольного набора переменных $V$ рассуждение можно
завершить ссылкой на лемму Цорна:%
\index{Лемма!Цорна}%
\glossary{\Цорн}
рассмотрим частично
упорядоченное множество, элементами которого будут
непротиворечивые множества формул, а порядком\т отношение \лк
быть подмножеством\пк. Рассуждение предыдущего абзаца
показывает, что всякая цепь в этом множестве имеет верхнюю
границу (объединение линейно упорядоченного по включению
семейства непротиворечивых множеств является непротиворечивым
множеством). Следовательно, для любого непротиворечивого
множества найдётся содержащее его максимальное непротиворечивое
множество. А оно обязано быть полным (иначе его можно расширить,
добавив $A$ или $\lnot A$).

Лемма~1 доказана.

Доказательство леммы~2. Пусть $\Gamma$\т непротиворечивое
полное множество. Тогда для каждой переменной (из множества~$V$) ровно
одна из формул $p$ и $\lnot p$ выводима из~$\Gamma$. Если
первая, будем считать переменную~$p$ истинной, если вторая\т
ложной. Тем самым появляется некоторый набор $\nu$ значений
переменных, и надо только проверить, что любая формула из
$\Gamma$ при таких значениях переменных истинна. Это делается
так: индукцией по построению формулы $A$ мы доказываем, что
        \begin{align*}
\text{$A$ истинна на наборе $\nu$}\ &\Rightarrow\ \Gamma\vdash A,\\
\text{$A$ ложна на наборе $\nu$}\ &\Rightarrow\ \Gamma\vdash
\lnot A.
        \end{align*}
Базис индукции (когда $A$\т переменная) обеспечивается определением
истинности переменных. Для шага индукции
используется та же лемма, что и при доказательстве
полноты с помощью разбора случаев. Пусть, например, $A$~имеет вид
$(B\land C)$. Тогда есть четыре возможности для истинности $B$ и
$C$. В одном из них (когда $B$ и $C$ истинны на $\nu$) по предположению
индукции мы имеем $\Gamma\vdash B$ и $\Gamma\vdash C$, откуда
$\Gamma\vdash (B\land C)$, то есть $\Gamma\vdash A$. В другом
($B$~истинна, $C$ ложна) предположение индукции даёт
$\Gamma\vdash B$ и $\Gamma\vdash \lnot C$, откуда
$\Gamma\vdash\lnot(B\land C)$, то есть $\Gamma\vdash\lnot A$.
Аналогично разбираются и все остальные случаи и логические
связки. Лемма 2 доказана, и тем самым завершено доказательство
теоремы~\ref{completeness-second-form}.
        %
\end{proof}

Мы доказали, что всякое непротиворечивое множество формул
совместно. Отсюда легко следует, что всякая тавтология
является теоремой. В самом деле, если $\varphi$\т тавтология,
множество $\{\lnot\varphi\}$ несовместно, поэтому из~$\lnot \varphi$
выводится противоречие, поэтому $\vdash \lnot\lnot\varphi$, и
по закону снятия двойного отрицания $\vdash\varphi$.

Кроме того, теорема о полноте во второй формулировке имеет такое
очевидное следствие:
        %
\begin{theorem}[теорема компактности для исчисления высказываний]
\index{Теорема!компактности для ИВ|defin}%
        \label{propositional-compactness}%
Пусть $\Gamma$\т множество формул, всякое конечное подмножество
которого совместно. Тогда и всё множество $\Gamma$ совместно.
        %
\end{theorem}

\begin{proof}
        %
Как мы знаем, несовместность равносильна противоречивости, а
вывод противоречия по определению может использовать лишь
конечное число формул.
        %
\end{proof}

Поскольку в формулировке теоремы компактности нет упоминания об
исчислении высказываний (речь идёт лишь об истинности формул, а
не о выводимости), возникает вопрос, нельзя ли её доказать
непосредственно.

\begin{problem}
        %
Дайте прямое доказательство теоремы компактности для случая,
когда переменных в множестве~$V$ конечное число. (Указание: любое несовместное множество имеет несовместное
подмножество мощности не больше $2^{|V|}$.)
        %
\end{problem}

Для случая счётного числа переменных можно воспользоваться
компактностью (в топологическом смысле слова) канторовского
пространства.%
\index{Канторовское пространство}
Его элементами являются бесконечные
последовательности нулей и единиц. Если две последовательности
отличаются в $n$\д й позиции, а все предыдущие члены совпадают,
то расстояние между ними считается равным~$2^{-n}$. Это
метрическое пространство компактно.

Пусть $V$ содержит счётное число переменных. Последовательность
значений переменных будем рассматривать как точку канторовского
пространства; формуле соответствует область, состоящая из точек,
где формула истинна. Поскольку формула содержит лишь конечное
число переменных, эта область является замкнутым и открытым
множеством одновременно. Пусть имеется множество формул, любое
конечное подмножество которого совместно. Это значит, что
соответствующие формулам подмножества канторовского пространства
образуют, как говорят, центрированную систему%
\index{Центрированная система множеств}
(любое конечное их
число имеет общую точку). А в компактном пространстве любое
центрированное семейство замкнутых множеств имеет общую точку
(иначе их дополнения образуют открытое покрытие, у которого нет
конечного подпокрытия). Эта их общая точка и будет набором
значений, на котором все формулы истинны.

То же самое рассуждение годится и для несчётного множества
переменных, но тогда возникает несчётное произведение
двухточечных пространств, которое является топологическим
пространством (но не метрическим); надо заметить, что это
пространство компактно по теореме~Тихонова,%
\index{Теорема!Тихонова}%
\glossary{\Тихонов}
после чего наше рассуждение проходит.

Для счётного набора переменных теорема компактности связана с
так называемой \emph{леммой Кёнига}.%
\glossary{\Кёниг}
Конечные последовательности
нулей и единиц (включая пустую последовательность) мы называем
двоичными словами. Двоичным деревом%
\index{Двоичное дерево}
мы называем множество
двоичных слов, которое вместе со всяким словом содержит все
его начала (начальные отрезки).
Бесконечной ветвью двоичного дерева $T$ мы называем
бесконечную последовательность нулей и единиц, любое конечное
начало которой принадлежит~$T$.

\begin{theorem}[лемма Кёнига]
\index{Лемма!Кёнига|defin}%
\glossary{\Кёниг}%
        %
Любое бесконечное дерево имеет бесконечную ветвь.
        %
\end{theorem}

\begin{proof}
        %
Говоря о бесконечности дерева, мы имеем в виду, что
соответствующее множество бесконечно. Отсюда следует, что оно
содержит слова сколь угодно большой длины. Пусть
$p_1,p_2,\dots$\т счётное множество переменных, которые
принимают значения $0$ или $1$. Для каждого $n$ рассмотрим
формулу $\varphi_n$, которая утверждает, что слово $p_1p_2\dots
p_n$ принадлежит дереву~$T$ (это возможно, так как любая булева
функция выразима формулой). Поскольку $T$\т дерево, $\varphi_i$
влечёт $\varphi_j$ при $j<i$. Любое конечное множество формул
вида $\varphi_i$ равносильно, таким образом, одной формуле с
максимальным~$i$ и потому совместно. Следовательно, и множество
всех формул $\varphi_i$ совместно, и выполняющий набор
определяет бесконечную~ветвь.
        %
\end{proof}

(Конечно, мы \лк бьём из пушек
по воробьям\пк: достаточно индукцией по~$i$ строить слово
длины~$i$, которое имеет бесконечное число продолжений в
дереве~$T$.)

Обычно утверждение леммы Кёнига формулируют так: если колония
бактерий, возникшая из одной бактерии, никогда не вымирает
полностью, то существует бесконечная последовательность
бактерий, каждая следующая из которых получается при делении
предыдущей. [Аналогичная формулировка про людей осложняется
возможностью клонирования, наличием двух полов и проблемами
политкорректности.]

\section[Поиск контрпримера и исчисление секвенций]
        {Поиск контрпримера и исчисление секвенций}

В этом разделе мы построим другой вариант исчисления
высказываний\т так называемое \emph{исчисление секвенций}.%
\index{Исчисление!секвенций|defin}
Такого рода исчисления изучаются в теории доказательств. Они
оказываются более удобными для анализа синтаксической структуры
выводов. Их называют исчислениями генценовского типа%
\index{Исчисление!генценовского типа}
(по имени
Генцена,%
\glossary{\Генцен}
который начал их изучать). Ранее приведённый вариант
исчисления высказываний называют исчислением гильбертовского
типа%
\index{Исчисление!гильбертовского типа}
(по имени Гильберта,%
\glossary{\Гильберт}
который использовал подобные
исчисления в своей программе формального построения математики).

Для начала мы вообще не будем говорить ничего об аксиомах и
правилах вывода, а рассмотрим задачу поиска контрпримера. Пусть
дана некоторая формула~$A$, про которую мы подозреваем, что она
не является тавтологией, и хотим найти значения переменных, при
которых она ложна. Как это делать? Естественно посмотреть на
структуру формулы~$A$. Например, если $A$ имеет вид~$B\to C$, то
надо найти значения переменных, при которых формула $B$ истинна,
а $C$ ложна. Если $A$ имеет вид~$B\lor C$, то надо найти
значения переменных, при которых обе формулы~$B$ и~$C$ ложны.
Если $A$ имеет вид~$B\land C$, то надо найти либо значения
переменных, при которых $B$~ложна, либо значения переменных, при
которых~$C$ ложна. Тем самым задача поиска контрпримера для
формулы~$A$ сводится к одной или нескольким задачам такого же или более
общего вида (надо обеспечить истинность или ложность одной или
нескольких формул).

Введём необходимую терминологию и обозначения. Будем называть
\emph{секвенцией}%
\index{Секвенция|defin}
выражение $\Gamma\vdash\Delta$, где $\Gamma$
и~$\Delta$\т некоторые конечные множества формул. (Пока что знак
$\vdash$ не имеет ничего общего с выводимостью, а только
разделяет два множества формул.) С каждой секвенцией $\Gamma\vdash\Delta$
будем связывать задачу поиска таких значений переменных, при
которых все формулы из~$\Gamma$ истинны, а все формулы из~$\Delta$
ложны. Такой набор значений мы по некоторым причинам будем
называть \emph{контрпримером}%
\index{Контрпример (к секвенции)|defin}
к секвенции $\Gamma\vdash\Delta$.
Легко проверить, что контрпримеры к секвенции $\Gamma\vdash\Delta$\т
это контрпримеры к формуле
       $$\land \Gamma\to\lor\Delta,$$
($\land\Gamma$ обозначает конъюнкцию формул из $\Gamma$,
а~$\lor\Delta$\т дизъюнкцию формул из $\Delta$), то есть те наборы
значений, при которых эта формула ложна. При этом конъюнкцию пустого
множества формул мы считаем тождественно истинной, а дизъюнкцию\т
тождественно ложной.

Наша исходная задача поиска контрпримера к формуле $A$ может
быть теперь сформулирована как задача поиска контрпримера к
секвенции $\vdash A$. (Мы позволяем себе писать так для
краткости; полностью следовало бы написать
$\varnothing\vdash\{A\}$.)

Задачу поиска контрпримера к
секвенции можно решать с помощью следующих правил. В каждом из
приведенных правил нижний заказ на контрпример выполним, если и
только если выполним один из верхних заказов, \те нижняя секвенция
имеет контрпример тогда и только тогда, когда хотя бы
одна из верхних секвенций имеет контрпример.
\begin{align*}
\frac{\Gamma\vdash A,\Delta\qquad \Gamma\vdash B,\Delta}
{\Gamma\vdash A\wedge B,\Delta}
        \qquad&\qquad
\frac{A,B,\Gamma\vdash\Delta}
{A\wedge B,\Gamma\vdash\Delta}
        \\[0.7ex]
\frac{\Gamma\vdash A, B,\Delta}
{\Gamma\vdash A\vee B,\Delta}
        \qquad&\qquad
\frac{A,\Gamma\vdash \Delta\qquad B,\Gamma\vdash \Delta}
{A\vee B,\Gamma\vdash \Delta}
        \\[0.7ex]
\frac{\Gamma,A\vdash  B,\Delta}
{\Gamma\vdash A\to B,\Delta}
        \qquad&\qquad
\frac{\Gamma\vdash A,\Delta\qquad \Gamma,B\vdash\Delta}
{A\to B,\Gamma\vdash \Delta}
        \\[0.7ex]
\frac{A,\Gamma\vdash \Delta}
{\Gamma\vdash \neg A,\Delta}
        \qquad&\qquad
\frac{\Gamma\vdash A,\Delta}
{\neg A,\Gamma\vdash \Delta}
\end{align*}

\noindent
Каждое из правил соответствует анализу одной из формул нижней
секвенции. Правила разделены на группы в зависимости от главной
связки анализируемой формулы, и согласованы с таблицами
истинности для этой связки (что легко проверить). Запятая в
правилах используется как сокращение: $\Gamma,A$ обозначает
$\Gamma\hm\cup\{A\}$ \итд

Как пользоваться этими правилами? Возьмём секвенцию,
к которой мы ищем контрпример. Выберем в ней одну
из формул слева или справа, посмотрим на главную
связку и применим соответствующее правило (написав
одну или две секвенции над исходной). Затем к каждой
из них снова применим одно из правил \итд Постепенно
будет расти \лк дерево поиска контрпримера\пк, причём
исходная секвенция будет иметь контрпример тогда и
только тогда, когда одна из верхних секвенций
(стоящих в \лк листьях\пк) этого дерева имеет
контрпример.

Когда этот процесс обрывается? Это происходит в
том случае, если все формулы в оставшихся секвенциях
представляют собой переменные, тогда ни одно из
наших правил поиска контрпримера не применимо. Но
к этому моменту всё становится ясным: если в левой
и правой части секвенции есть общая переменная,
то к ней нет контрпримера (одна и та же переменная
не может быть одновременно истинной и ложной). Если
же левая и правая часть такой секвенции не
пересекаются, то контрпример есть.

Вот как это делается с секвенцией ${\vdash(p\to q)}\hm\to q$:
        $$
\frac{\mathstrut\vdash p,q\;\qquad q\vdash q}
{\frac{\displaystyle \mathstrut(p\to q)\vdash q}
{\displaystyle \mathstrut\vdash (p\to q)\to q}}
         $$
Контрпример найден: $p$ и $q$ ложны (он является контрпримером к
секвенции $\vdash p,q$).

Напротив,
поиск контрпримера к секвенции
${\vdash((p\to q)\to p)}\hm\to p$
не даёт результата:
       $$
\frac{\frac{%
\frac{\displaystyle \mathstrut p\vdash p,q}
     {\displaystyle \mathstrut\vdash p,p\to q}\qquad
\displaystyle
     \mathstrut p\vdash p}%
{\displaystyle \mathstrut (p\to q)\to p\vdash p}}%
{\mathstrut\vdash((p\to q)\to p)\to p}
        $$
Здесь обе секвенции $p\hm\vdash p,q$ и $p\hm\vdash p$ не имеют
контрпримеров. Следовательно, формула $((p\to q)\to p)\to p$
является тавтологией.

Построенный алгоритм можно одновременно рассматривать
как доказательство полноты некоторого \лк исчисления
секвенций\пк.

\emph{Аксиомами}%
\index{Аксиомы!исчисления секвенций|defin}%
\index{Исчисление!секвенций!аксиомы|defin}
исчисления секвенций будем называть секвенции,
в левых и правых частях которых встречаются только переменные,
причём некоторая переменная встречается в обеих частях.

\emph{Правилами вывода}%
\index{Правила вывода исчисления секвенций|defin}%
\index{Исчисление!секвенций!правила вывода|defin}
в исчислении секвенций являются правила
нашей таблицы. Каждое из этих правил объявляет выводимой нижнюю
секвенцию, если выводимы все верхние. (Процесс вывода
естественно представлять в виде дерева, как в наших примерах, но
можно развернуть и в последовательность секвенций.)

\begin{theorem}[корректность и полнота исчисления секвенций]
\index{Теорема!о корректности и полноте исчисления секвенций|defin}%
        %
Секвенция выводима тогда и только тогда, когда
она не имеет контрпримера.
        %
\end{theorem}

\begin{proof}
        %
Аксиомы не имеют контрпримера. Если все верхние секвенции какого\д то
правила вывода не имеют контрпримера, то и нижняя секвенция не
имеет контрпримера. (Именно так мы подбирали правила: контрпример
к нижней секвенции будет контрпримером к одной из верхних.)
Следовательно, все выводимые секвенции не имеют контрпримера.

Обратно, пусть секвенция не имеет контрпримера. Тогда
описанный процесс поиска контрпримера обрывается на
аксиомах и тем самым даёт её вывод.
        %
\end{proof}

В частности, если формула $A$ является тавтологией, то секвенция
${\vdash A}$ выводима в исчислении секвенций. Это обстоятельство
можно использовать для ещё одного доказательства полноты
исчисления высказываний (в стандартной, гильбертовской форме). В
самом деле, для каждой секвенции $\Gamma\vdash\Delta$ рассмотрим
\emph{представляющую её формулу},%
\index{Формула!представляющая секвенцию|defin}%
\index{Секвенция!представляющая формула|defin}
то есть формулу ${\land
\Gamma} \hm\to {\lor\Delta}$ (в левой и правой частях стоят
конъюнкция и дизъюнкция формул из $\Gamma$ и $\Delta$
соответственно). Теперь индукцией по построению вывода в
исчислении секвенций надо доказать такой факт: если секвенция
выводима в исчислении секвенций, то представляющая её формула
выводима в (обычном) исчислении высказываний.

Это требует довольно хлопотной проверки, впрочем. Сначала надо
проверить, что конъюнкция и дизъюнкция доказуемо ассоциативны, и
потому всё равно, как расставлять скобки в формуле,
представляющей секвенцию. Затем полезно убедиться, что правила,
связанные с отрицанием (два последних правила таблицы) можно
применять в обе стороны, не меняя выводимости соответствующей
формулы. Другими словами, надо проверить, что формулы
        $$
 \Gamma\to(\Delta\lor A) \quad \text{и} \quad
 (\Gamma\land\lnot A)\to\Delta,
        $$
а также формулы
        $$
 \Gamma\to(\Delta\lor \lnot A) \quad \text{и} \quad
 (\Gamma\land A)\to\Delta
        $$
выводимы одновременно (здесь $\Gamma$ и $\Delta$ можно считать
формулами, а не множествами формул, расставив скобки надлежащим
образом). После этого мы можем переносить формулы из одной части
в другую (меняя их на противоположные), и потому можем считать
одно из множеств $\Gamma$ и $\Delta$ пустым, если это нам
удобно. Теперь ссылка на лемму о дедукции и уже известные нам
свойства выводимости завершает рассуждение.

\begin{problem}
        %
Провести это рассуждение подробно.
        %
\end{problem}

Естественно возникает вопрос\т чем уж так интересно исчисление
секвенций? Какая, собственно говоря, разница\т иметь дело с
секвенциями или с формулами, раз всякую секвенцию можно
представить формулой? Принципиальное различие тут в следующем.
Правила вывода в исчислении секвенций таковы, что в их верхнюю
часть входят только подформулы формул, встречающихся в нижней
части. Поэтому в выводе какой\д то секвенции не может встретиться
ничего принципиально нового, чего не было в самой секвенции. В
гильбертовском исчислении это далеко не так: мы можем вывести
формулу $B$ из формул $A\to B$ и $A$, при этом формула $A$ может
быть совершенно произвольной. Это же различие объясняет, почему
поиск вывода снизу вверх (как можно теперь называть то, что
раньше называлось поиском контрпримера\т мы находим либо
контрпример, либо вывод) для исчисления секвенций происходит
сравнительно однозначно (мы можем по\д разному выбирать
расчленяемую формулу, но и только), в то время как искать вывод
в обычном исчислении высказываний, начав с интересующей нас
формулы $B$ и смотря, из чего бы она могла получиться, не
удаётся (если только не перебирать все формулы подряд).

Заметим, что добавление к исчислению секвенций уже упоминавшегося
правила сечения
        $$
\frac{\Gamma\vdash \Delta,A \qquad \Gamma, A\vdash \Delta}
     {\Gamma \vdash \Delta}
        $$
нарушает свойство \лк подформульности\пк, так как формула~$A$
может быть никак не связана с нижней секвенцией. Отметим, что
добавление правила сечения не нарушает корректности, как легко
проверить, и не может нарушить полноты.

Задачи о поиске вывода и анализе его структуры, хотя и были одно
время модными в связи с \лк искусственным интеллектом\пк,
представляют большой интерес, и про них есть целая наука, в
которой исчисления генценовского типа играют центральную роль.
Мы рассмотрели один из вариантов исчисления секвенций для
классического исчисления высказываний; бывают исчисления для
интуиционистских и модальных логик, для исчисления предикатов
\итд Исходной мотивацией для рассмотрения такого рода исчислений
было желание доказать непротиворечивость арифметики. Согласно
знаменитой
второй теореме Гёделя%
\glossary{\Гёдель}
о неполноте%
\index{Теорема!Гёделя о неполноте (вторая)}
без дополнительных аксиом
этого сделать нельзя, но если принять схему аксиом для
трансфинитной индукции по ординалу $\varepsilon_0$, это удаётся
сделать, как показал Генцен.%
\glossary{\Генцен}
